\section{公诉案件第一审程序}
\begin{enumerate}
  \item 检察院。
        \begin{enumerate}
          \item \zx{审理中新的事实}。新 = 影响定罪量刑 \hz 需要补查补证。
                \begin{enumerate}
                  \item 漏罪 \yjt \zx{补充}起诉。
                  \item 漏人 \yjt \zx{追加}起诉。
                  \item \zx{法条/身份错 \yjt 变更}起诉。\sjt \al{可口头进行 \dui}。
                  \item 没查清 \yjt 补充侦查。
                \end{enumerate}
          \item \zx{无罪无责方可撤诉}。
        \end{enumerate}
  \item 法院。
        \begin{enumerate}
          \item 发现新的事实应当\zd{通知}检。\al{撤回 \cuo}。
          \item 分案、并案具体问题具体分析。\zd{做题，提到分案基本都对}。
          \item 庭前审查。决定\zd{能否开庭}审判。各事项是否就位。\zd{7 日内处理，计入审期}。
                \begin{enumerate}
                  \item \zd{必经程序}。\zx{检提起公诉的案件，都应当受理}。书面，程序性审查（都审了后面干啥）。
                        \al{因被害人不断上访重新起诉的，不予受理 \cuo}。
                  \item \zx{缺材料 \yjt 3 日内补送}。可以是各种材料。\al{发现有\xhx{自首情节}，移送\xhx{材料中没有}}。
                  \item \zx{身份不明}正常受理、审理。调查核实 \hz 用其自报的姓名。\bs{身份不明不按缺材料的情况进行处理}。
                  \item \bs{退回。管辖错；不在案（缺席审判除外）；无新证，再起诉；过时效，特赦告诉和死掉}。
                  \item \bs{告诉才处理的 \yjt 退回 + 告知被害人提起自诉}。
                \end{enumerate}
          \item 庭前准备。\bs{通知、送达、Opt(公告)}。
                \begin{enumerate}
                  \item 确定审判长和合议庭人员。\yi 随机。\er 书记员不属于合议庭成员。
                        \al{甲为一审书记员，案件被发回重审，甲不能再担任本案书记员 \cuo}。
                  \item \zd{10} 日前。起诉书副本送达\zd{被告人、辩护人}。\zd{送晚了不算程序违法} \yjt 申请变更开庭时间。
                  \item \zd{5} 日前。提供\zd{证}人、鉴定人名单，以及拟当庭出示的证据。
                  \item \zd{3} 日前。通知检；传票和通知书等送达\zd{当事人}；公开的，公布案由、被告人姓名、开庭时间地点。
                \end{enumerate}
          \item Opt(庭前会议)。一般不公开。\bs{程序性审查，而非实质性审查}。
                \begin{enumerate}
                  \item \yi 依申请 \yjt 法院决定。\er 依职权。\zx{多、大、争议}。材料多、案情重大复杂、事实和证据争议较大。
                  \item \zx{不是被告人权利，能否参加看法院}。\al{应通知甲参加 \cuo}。
                  \item 内容。出庭人员名单；排非、重新鉴定/勘验；\bs{管辖权异议}、回避、是否公开。
                  \item \zx{明显事实不清，证据不足}的案件 \yjt \zx{建议撤诉}（\al{退回 \cuo}）。
                  \item \bs{无异议}证据庭审时举证、质证\bs{可简化}（\al{不出示 \cuo}）。If(庭中反悔) + 正当理由才处理。
                  \item 可能导致庭审中断的\bs{程序性事项，可在会议后处理，并在庭审中说明处理决定和理由}。\al{如回避}。
                  \item \al{申请排非的，可在庭前会议中就是否排非\xhx{作出决定} \cuo}。
                  \item \al{可以讨论\xhx{并确定}罪名}。
                \end{enumerate}
          \item 开庭环节。书记员先进，查明人员是否到场。审判长，查明基本信息。
          \item \zx{基本流程}。\zx{起诉书} \yjt 宣布\zx{排非}情况 \yjt 讯问被告人 \yjt 询问证人、鉴定人 \yjt 出示物证 \yjt 宣读\zx{勘验}笔录 \yjt \zx{公诉词} \yjt 最后陈述。
          \item 法庭调查。\bs{主要是事实 + 证据}。
                \begin{enumerate}
                  \item 公诉人宣读起诉书。Exist(口头、书面实质性差异) \yjt 指定时间内书面方式提出。书面为准。
                  \item 被告人、被害人陈述。
                  \item 讯问（公权力的人）被告人（不是必须、多名分别进行），\zd{发问}被害人。
                  \item 出示、核实证据。\zd{先人后物再笔录}。\zx{申请方先问}；\zd{分别进行，不得旁听}。
                  \item \zx{允许证据突袭 \hz 调取新证据}。有理由 + 有必要 + 法院同意 \yjt 双方进行质证。同民诉。
                        \al{庭前没有出示的证据，庭审\xhx{不得}出示 \cuo}。
                  \item 补充侦查。
                  \item 合议庭调查核实证据。% \zd{唯一没有搜查}。不影响，非关键，对被告人有利的证据。
                  \item 对量刑事实、证据的调查。
                  \item 对查扣冻财物的调查。
                  \item \zx{证据问题，都得听听在做各位的意见}。
                \end{enumerate}
          \item 法庭辩论。\zd{主要是法律适用}。\bs{先控方后辩方，先定罪后量刑}。
                \begin{enumerate}
                  \item 公诉人 \yjt 被害人及其诉讼代理人 \yjt 被告人自行辩护 \yjt 辩护人辩护 \yjt 控辩双方进行辩论。
                \end{enumerate}
          \item 最后陈述。\zd{不能省略}。可以制止；应当制止；程序可能回转。
          \item \zd{评议、宣判}。当庭 5 日内送达；定期立即送达。
                \begin{enumerate}
                  \item 有罪判决。指控 A 罪 \yjt 认定 B 罪 \yjt 判决 B 罪。\zx{指控的这个事实判成啥罪，法院说了算}。
                        \al{可以直接\xhx{自行变更}罪名 \dui}。
                        \zd{应当在判决前听取双方意见，保障被告人辩护权，必要时可再开庭，组织控辩双方围绕构成何罪进行辩论}。
                  \item \zx{无罪判决}。
                        \begin{enumerate}
                          \item 证据不足（疑罪从无）。\zx{发现新证据 \yjt 重新起诉}。不能撤销原判，新判决书写明曾判无罪
                                \footnote{再审是针对错误的判决，这种情况你不能说他判错了，既然没错，那就不用撤。}。
                          \item 部分有罪，部分无罪。事实\zx{不清部分，不予认定}。
                        \end{enumerate}
                  \item 判决不负刑事责任。干了但\zx{无责任}年龄 \hz 无责任能力。\zx{未成年 16、精神病}。
                        \al{经审理发现被告系\xhx{精神病人}，在不能控制自己行为时造成危害结果}。
                \end{enumerate}
          \item \zd{审期}。2 + 1（院长） + 3（交、集、流、广、死、附民 \yjt 上一级法院）+ X（最高法）。
          \item \zx{审判障碍}。
                \begin{enumerate}
                  \item 休庭。\al{辩护人被带出法庭，应当决定延期审理 \cuo}。
                  \item \zx{延期}审理。用决定决定决定。这次开不了，换个时间再开。\bs{新证、补侦、和回避}。
                  \item \bs{中止}审理。用裁定裁定裁定。\bs{他跑、他病耐心等}。\al{新冠算不可抗力 \dui}。
                  \item \bs{终止审理。过时效，特赦告诉和死掉}。\al{发现是侵占罪}。
                \end{enumerate}
          \item 法庭秩序。
                \begin{enumerate}
                  \item 警告、训诫、责令退出法庭、强行带出法庭；\zx{1000 元以下} 罚款 \hz \zx{15} 日以下拘留。\zd{院长决定}。
                        \al{可以先行拘留 \cuo}。
                  \item 违规使用录音录像设备\zx{暂扣}，删除后归还。
                \end{enumerate}
          \item 判决书。
                \begin{enumerate}
                  \item 首部。标题、编号、公诉机关和诉讼参加人情况、案件审理的有关情况。
                  \item 事实部分。指控的基本内容、辩护意见、查明的事实和据以定罪的证据。
                  \item \zx{理由部分。罪名、理由、法律适用}。
                  \item 尾部。交代上诉权和上诉法院、合议庭组成人员署名、判决的决定日期、书记员署名。
                \end{enumerate}
          \item \bs{开庭问题}。
                \begin{enumerate}
                  \item 一审、事实、证据。一审 = 一审、[二审、再审发回用一审程序审的]。
                  \item 一审有抗诉、二审要他死、再审 +重罚。“核”准类程序。\al{上诉 \cuo}。
                  \item 减刑假释。\zd{重大立功不一般；重大影响有意见；有权有势又有钱。职务黑金}。
                  \item 没收程序。一般不开庭，利害关系人\zd{到庭参诉}。\zx{审钱不审人当然不开庭}。
                  \item 强疗。一般要开庭，法代请求不开可以不开。\zx{审人当然开庭}。
                  \item \zx{开庭、同级检必派员 \yjt 不开庭不派员}。
                \end{enumerate}
          \item \zd{期限总结}。
                \begin{enumerate}
                  \item 侦查羁押。2 + 1 + 2 + 2 + X。
                  \item 审查起诉。1 + 15 日。
                  \item 一审。2 + 1 + 3 + X。
                  \item 二审。2 + 2 + X。
                \end{enumerate}
          \item 发问总结。\zx{谁申请谁先问}。
                \begin{enumerate}
                  \item 针对。证人、鉴定人、有专门知识的人的出庭。
                  \item 向法庭作陈述 \yjt 审判长许可 \yjt 谁申请谁先问 \yjt 另一方问。
                \end{enumerate}
        \end{enumerate}
\end{enumerate}
